\chapter{Literature Review and Related Work}
\label{ch:literature}

This chapter examines the research needed to justify Manzil's design rather than treating the use of familiar tools as a contribution. The review is organized around five decisions: how traffic is observed, how normal signals are controlled, how learned policies are evaluated, how emergency routes are selected, and how signal priority is activated and recovered. Each study is considered in its original data or simulation setting because reported values from incompatible benchmarks cannot be treated as a common comparison. The discussion identifies the method, evidence, limitations, and direct implication for a Manzil design choice. Research papers are distinguished from official software documentation, which is used to establish interface capabilities but is not counted toward the study requirement. The purpose is not to assemble every publication on intelligent transportation. It is to establish whether the proposed components are technically defensible, which assumptions must be recorded, and where a controlled undergraduate experiment can add useful evidence. The final synthesis consequently states a precise integration question instead of claiming that the use of YOLO, reinforcement learning, or shortest-path search is itself novel.

\section{Review Method and Terminology}

The approved proposal supplied the initial bibliography, after which each retained study was checked against a publisher page, proceedings record, DOI, or official open-access paper. Selection required direct relevance to vehicle detection and tracking, traffic-state estimation, signal control, simulation, graph routing, or emergency preemption. The review also applies consistent terminology. Detection refers to locating and classifying vehicles within one frame, while tracking refers to associating those observations through time. A detector therefore cannot produce reliable passage counts by itself. Signal priority describes preferential treatment that may remain within normal timing constraints, whereas preemption may interrupt normal operation for an emergency request. Manzil uses the term moving green corridor for a coordinated sequence of bounded interventions ahead of the ambulance, not for a permanently green route. These distinctions affect the interpretation of prior results and the design of the proposed experiment. A paper reporting benchmark detection accuracy does not establish a calibrated queue estimate, and a study of isolated preemption does not automatically validate network-level routing. The following sections preserve those boundaries when relating the literature to Manzil.

\section{Vision-Based Traffic Observation}

\subsection{One-Stage Object Detection}

Redmon et al. formulated YOLO as a single network that predicts bounding boxes and class probabilities from an image in one evaluation \cite{redmon2016yolo}. The study established the feasibility of real-time one-stage detection and reported a strong speed advantage, while also identifying localization error as a relative weakness. The design is relevant to Manzil because a compact inference path is practical for recorded traffic video and may later support a live demonstration on available hardware. The evidence does not establish accuracy for every vehicle class, road layout, camera height, weather condition, or Lahore traffic scene. A selected pretrained checkpoint must therefore be evaluated on the project's own held-out sequences, and confidence or intersection-over-union thresholds cannot be copied from an unrelated benchmark. Frame-level boxes also do not provide stable counts by themselves. The same vehicle may be detected in many frames, an occluded vehicle may disappear temporarily, and perspective changes apparent size across the image. Manzil consequently treats one-stage detection as the first part of a measurement pipeline that also requires tracking, road-region assignment, temporal aggregation, and explicit uncertainty.

\subsection{Multi-Object Association}

Wojke et al. extended a simple online tracker with a learned appearance metric to improve association through longer occlusions \cite{wojke2017deepsort}. Appearance information can reduce identity switches, but it adds a feature-extraction model, processing cost, and sensitivity to domain differences. ByteTrack follows a different strategy by associating high-confidence detections first and then attempting to recover true objects from lower-confidence boxes rather than discarding them \cite{zhang2022bytetrack}. The authors evaluated it across several tracking benchmarks and reported strong accuracy with real-time throughput on their stated hardware. This approach is relevant to traffic video because occlusion and motion blur often reduce detection confidence. Neither benchmark establishes Manzil's processing speed, lane assignment, or queue accuracy. An identity-preserving track still requires a road-contact point, direction rule, and region geometry before it becomes a valid traffic observation. These studies therefore motivate an empirical tracker choice based on held-out count error, identity stability where identities are needed, failure under occlusion, and end-to-end throughput on the team's hardware. If a simpler association method produces equally reliable aggregate counts, its lower complexity may be preferable.

\subsection{Traffic Datasets and Domain Coverage}

BDD100K provides a large and diverse driving dataset with detection, tracking, and other driving-scene annotations collected across varied locations, weather, and times \cite{yu2020bdd100k}. Its diversity is useful for model initialization and stress testing, although its primarily vehicle-mounted viewpoint differs from a fixed roadside or elevated camera. UA-DETRAC focuses directly on vehicle detection and tracking in traffic-surveillance video, making its camera geometry and evaluation protocol closer to an intersection-observation task \cite{wen2020uadetrac}. Its scenes still cannot be assumed to represent Lahore traffic composition, motorcycle and rickshaw prevalence, lane discipline, lighting, or camera placement. The two datasets are therefore complementary references rather than automatic substitutes for local validation. Manzil should use licensed public sequences where they support reproducible evaluation and may add a small, separately documented local sample if collection permission and annotation effort are feasible. These datasets also influence the split strategy. Adjacent video frames are highly correlated, so a random frame split can place nearly identical observations in training and test data. The project will split by source or continuous sequence, retain a held-out evaluation set, and avoid tuning thresholds on that set.

\section{Conventional and Adaptive Signal Control}

\subsection{Fixed-Time Baseline}

A fixed-time controller remains necessary even when it is not expected to dominate every scenario. It provides a stable reference against which the value of traffic-responsive decisions can be measured, and it helps reveal whether demand generation or metric collection is biased before more complex controllers are introduced. The baseline must be credible rather than deliberately weak. Manzil will therefore document its phase sequence, green times, yellow and all-red clearance, cycle length, offsets where relevant, minimum service, and the demand conditions used to choose those settings. The same phase definitions and safety transitions will be shared with adaptive controllers so that a comparison does not reward one controller for using an illegal or more permissive movement set. Fixed timing is also useful for route and corridor ablations because its behavior is deterministic and easy to reproduce. If the fixed plan is severely under-saturated or over-saturated in every scenario, differences in controller performance will be difficult to interpret. A basic capacity and queue check will therefore confirm that the scenario family contains meaningful light, medium, and heavy conditions instead of a single demand level chosen to favor adaptation.

\subsection{Max-Pressure Control}

Varaiya formalized max-pressure control for a network of signalized intersections using differences between upstream and downstream queue conditions \cite{varaiya2013maxpressure}. Selecting a phase that reduces network pressure has a clearer transportation basis than combining arbitrary percentages without explaining what the score represents. The approach is responsive and can operate from local movement information, which makes it a valuable non-learning reference for Manzil. Its assumptions still matter. Queue observations must correspond to valid movements, downstream storage must be represented, and a theoretical selection rule does not remove practical minimum-green, yellow, all-red, or maximum-red constraints. Manzil's initial rule-based adaptive controller will therefore use a pressure-inspired score only when the required upstream and downstream measures are available and sufficiently fresh. Waiting time or maximum-red information will remain as an explicit fairness safeguard so that a low-volume approach is not postponed indefinitely. The controller will also provide a reason for each request, allowing reviewers to see which queues and constraints influenced the decision. This transparency makes the baseline useful even if the later reinforcement-learning policy does not outperform it.

\section{Reinforcement Learning for Traffic Signals}

\subsection{Cycle-Level Deep Q-Learning}

Liang et al. used deep reinforcement learning to adjust traffic-light cycle behavior from observed traffic state \cite{liang2019dqn}. The study demonstrates how an action definition, decision interval, state representation, and reward can be joined into a trainable signal controller. It also illustrates a limitation common to traffic RL: performance obtained for a particular intersection, simulator configuration, demand pattern, and training budget cannot be assumed to transfer to a different network. Manzil will therefore not reproduce a reported improvement as an expected project target or select deep Q-learning merely because it is familiar. A value-based method is feasible only after the discrete action space is bounded, legal phase transitions are enforced outside the learned policy, and fixed and rule-based baselines behave correctly. State normalization and reward components must be declared before evaluation, while training and tuning scenarios must remain separate from held-out tests. The policy will be assessed using traffic measures and safety interventions, not training return alone. If a simpler controller performs as well under the selected scenarios, that result will be reported rather than interpreted as a failure to produce a sufficiently complex model.

\subsection{Pressure-Informed Reward Design}

PressLight connects reinforcement-learning state and reward design to the transportation concept of max pressure \cite{wei2019presslight}. The paper argues that a theoretically grounded pressure signal can reduce reliance on arbitrary handcrafted objectives and evaluates the method against conventional and learning baselines in arterial networks. This is relevant to Manzil because a reward based only on total waiting time may conceal spillback, unbalanced approaches, or repeated postponement of a minor movement. Pressure can provide a principal network-efficiency signal, while explicit penalties or constraints can address starvation and unnecessary switching. The study does not remove the need for careful evaluation. Topology, movement mapping, flows, phase definitions, observation range, and training budget all influence the learned result. Manzil must therefore record the exact state features, normalization, action timing, reward components, and safety interventions used by the selected policy. Emergency priority will initially remain a supervisory state machine rather than a rare reward event folded into normal training. That separation permits the normal-controller comparison to be interpreted independently and lets the additional effect of the corridor be measured through a controlled ablation.

\subsection{Phase Competition and Network Cooperation}

FRAP represents competition among signal phases and exploits structural relationships so that equivalent movements need not be learned as unrelated cases \cite{zheng2019frap}. This is useful where an intersection contains several competing phases, but it requires an accurate phase graph and a more specialized representation. CoLight addresses cooperation across intersections through graph attention, allowing neighboring conditions to influence an agent's decision \cite{wei2019colight}. Chu et al. similarly study decentralized multi-agent deep reinforcement learning at larger scale and emphasize the importance of communication and partial observability \cite{chu2020marl}. Together, these studies show why the phrase ``multi-intersection RL'' is not a complete method. A project must state whether one centralized policy chooses all actions, independent agents operate without communication, or cooperating agents exchange defined observations. Multi-agent learning also introduces non-stationarity because each agent experiences an environment affected by policies that are changing at the same time. For Manzil, attention-based coordination may exceed the available tuning budget and make attribution difficult. The initial design will therefore use a small network and one explicit ownership model, adopting useful phase-structure ideas without claiming to reproduce every cited architecture.

\subsection{Simulation Environments}

CityFlow was introduced as a reinforcement-learning environment designed for large-scale city traffic and high simulation throughput \cite{ma2019cityflow}. It is an important alternative when training speed and very large networks dominate the design. SUMO is retained for Manzil because the approved proposal already depends on its microscopic vehicle model, traffic-light programs, OpenStreetMap conversion tools, NetEdit workflow, and TraCI interface. Lopez et al. describe the simulator's architecture and its support for detailed, externally controlled experiments \cite{lopez2018sumo}. This combination is appropriate for a manageable multi-intersection prototype in which route validity, individual ambulance movement, signal transitions, and per-vehicle outcomes must be inspected. Retaining SUMO also avoids the additional work of showing that a migrated network and demand model behave equivalently in another simulator. The choice does not make simulation results equivalent to field evidence. Car-following parameters, lane connections, signal programs, vehicle insertion, and map cleanup can materially change outcomes. Manzil will therefore version the network, demand, routes, simulator settings, and software release, and it will limit conclusions to the declared simulation conditions.

\section{Dynamic Routing}

Dijkstra's algorithm finds a minimum-cost path on a graph whose edge weights are nonnegative \cite{dijkstra1959}. Hart, Nilsson, and Raphael established the formal basis of A*, which adds a heuristic lower-bound estimate to guide the search \cite{hart1968astar}. Both methods are well established, so selecting either algorithm is not a research contribution by itself. Manzil's project-specific task is to build a graph that respects valid SUMO vehicle movements and to define a cost that reflects time rather than distance alone. The proposed cost combines link travel time with non-overlapping estimates of queue discharge and residual signal delay, all expressed in seconds. If measured link travel time already includes queue delay, a separate congestion term cannot be added without double-counting. Dijkstra will provide the reference route, while A* is permitted only when free-flow travel time forms an admissible lower bound under nonnegative added delay. Dynamic rerouting also needs stability. A route will change when the current path becomes infeasible or when an alternative offers a persistent improvement above a defined threshold. Route revisions, no-route cases, and computation time will be evaluated rather than treated as incidental log messages.

\section{Emergency Routing and Signal Preemption}

Humagain et al. systematically reviewed emergency-vehicle route optimization and preemption methods and identified recurring concerns involving dynamic traffic data, travel-time objectives, ordinary-traffic effects, safety, and limited real-world application \cite{humagain2020review}. The review supports an integrated evaluation but does not justify the unsupported claim that integrated solutions do not exist. Qin and Khan examine transitions into emergency preemption and recovery to normal timing, including clearance and the queues created by interruption \cite{qin2012preemption}. Their work shows that preemption is a lifecycle rather than a single instruction to set an approach green. Louati et al. investigate queue-aware emergency crossing across signalized intersections \cite{el2018emergency}. This is relevant because a green signal does not guarantee an unobstructed approach when the standing queue cannot discharge before the emergency vehicle arrives. Manzil therefore models requested, safe-transition, active-priority, passage, recovery, and cancellation states. It also uses look-ahead activation, timeout, downstream checks, and bounded priority duration. The precise timing values remain configurable simulation policies, not certified real-world safety parameters, and the evaluation will report ordinary-traffic disruption together with ambulance benefit.

\section{Synthesis and Research Gap}

The literature establishes reliable building blocks while exposing the risks that arise when they are connected. One-stage detection is fast enough to investigate traffic video, but useful traffic indicators require tracking, road geometry, temporal rules, and validation on the intended viewpoint. Max-pressure and reinforcement-learning research provide responsive control strategies, yet their comparison depends on state representation, phase constraints, reward design, training data, and simulator configuration. Dijkstra and A* solve well-defined path problems, but emergency routing remains only as credible as its movement graph, time-based costs, observation freshness, and rerouting stability. Preemption research makes safe transition and recovery central, and the review of emergency methods shows why ordinary-traffic disruption cannot be omitted. Manzil's proposed contribution is consequently not a new detector, shortest-path algorithm, or reinforcement-learning family. It is the controlled evaluation of an integrated and traceable prototype in which perception and simulation evidence remain distinct, normal controllers receive matched scenarios, routing and corridor effects are separated, and ambulance benefit is reported with network disruption and recovery. The unresolved question is whether that integration yields consistent benefits under the selected scenarios and constraints.

Table \ref{tab:literature-summary} consolidates the evidence discussed in this chapter so that the basis for each design choice can be checked without treating unlike experiments as a numerical league table. The second column records the method and the setting in which evidence was obtained, while the remaining columns separate limitations from the decision that Manzil draws from the study. This distinction is important because benchmark accuracy, simulator throughput, and travel-time improvements measure different constructs and cannot be ranked on one common scale. The table also makes the review auditable: each row is tied to a cited source, each limitation identifies a boundary that the project must respect, and each implication points forward to a requirement, design element, or evaluation control. Portrait presentation is retained by combining the method and evidence into one wrapped column instead of rotating the page. The resulting summary is deliberately concise; detailed qualifications remain in the preceding thematic discussion, where datasets, assumptions, and transfer limitations can be explained in context.

{\small
\setlength{\tabcolsep}{3pt}
\begin{longtable}{p{0.15\textwidth}p{0.29\textwidth}p{0.24\textwidth}p{0.22\textwidth}}
\caption{Literature review summary}\label{tab:literature-summary}\\
\toprule
\textbf{Study} & \textbf{Method and evidence} & \textbf{Limitations for Manzil} & \textbf{Design implication}\\
\midrule\endfirsthead
\toprule
\textbf{Study} & \textbf{Method and evidence} & \textbf{Limitations for Manzil} & \textbf{Design implication}\\
\midrule\endhead
Redmon et al. \cite{redmon2016yolo} & Unified one-stage detection; demonstrated real-time operation and a useful speed--accuracy tradeoff on standard benchmarks. & Original classes and benchmarks do not establish roadside traffic accuracy. & Evaluate a pretrained descendant on selected traffic video.\\
Wojke et al. \cite{wojke2017deepsort} & Online tracking with learned appearance association; improved association through occlusion in its pedestrian setting. & Adds computation and a domain-dependent appearance model. & Consider it when identity fragmentation affects counts.\\
Zhang et al. \cite{zhang2022bytetrack} & Two-stage association that retains low-score boxes; strong results on several tracking benchmarks. & Benchmark identities and speed do not imply calibrated traffic measures. & Measure end-to-end count error before selecting the tracker.\\
Yu et al. \cite{yu2020bdd100k} & Large multi-task driving-video dataset with varied conditions and annotations. & Mostly vehicle-mounted viewpoints differ from roadside Lahore scenes. & Use for initialization or stress tests, not as sole validation.\\
Wen et al. \cite{wen2020uadetrac} & Vehicle detection and tracking benchmark for surveillance scenes with a joint evaluation protocol. & Vehicle mix and site geometry do not represent every local scene. & Split by sequence and report local-domain limitations.\\
Varaiya \cite{varaiya2013maxpressure} & Max-pressure network signal control with theoretical properties under stated model assumptions. & Requires valid movement queues, downstream state, and lawful phases. & Use pressure as a transparent adaptive reference.\\
Liang et al. \cite{liang2019dqn} & Deep Q-learning for cycle control; demonstrated learned adaptation in a simulated setting. & Transfer across demand patterns and topology remains uncertain. & Bound actions and evaluate on held-out scenarios.\\
Wei et al. \cite{wei2019presslight} & Reinforcement learning with pressure-grounded state and reward; exceeded tested baselines. & Findings depend on the reported networks, flows, and training setup. & Prefer a grounded reward to an arbitrary weighted sum.\\
Zheng et al. \cite{zheng2019frap} & Phase-competition representation that learns relationships among competing movements. & Requires a more specialized phase encoding. & Preserve phase legality and structural invariance where feasible.\\
Wei et al. \cite{wei2019colight} & Graph-attention cooperation across signals; demonstrated coordination in tested simulations. & Training and attribution are more complex. & Keep the first network small and declare the coordination model.\\
Chu et al. \cite{chu2020marl} & Multi-agent deep reinforcement learning for larger networks and decentralized coordination. & Non-stationarity and tuning cost complicate reproducible comparison. & State ownership and communication explicitly.\\
Ma et al. \cite{ma2019cityflow} & High-throughput city-scale traffic-learning simulator intended for large experiments. & Migration would complicate alignment with the approved proposal. & Retain SUMO and document CityFlow as an alternative.\\
Lopez et al. \cite{lopez2018sumo} & Microscopic simulator architecture supporting detailed external control. & Network and car-following parameters limit external validity. & Version the network, demand, and simulator settings.\\
Dijkstra \cite{dijkstra1959} & Nonnegative shortest-path search with a minimum-cost guarantee under its assumptions. & Represents a static snapshot unless costs are refreshed. & Use as the reference route implementation.\\
Hart et al. \cite{hart1968astar} & Heuristic minimum-cost search that guides exploration when the heuristic is valid. & An inadmissible heuristic can lose the optimality guarantee. & Use a free-flow-time lower bound or retain Dijkstra.\\
Humagain et al. \cite{humagain2020review} & Systematic review of emergency routing and preemption research. & A review does not validate any one implementation. & Measure emergency benefit and ordinary-traffic impact.\\
Qin and Khan \cite{qin2012preemption} & Simulation of preemption entry and recovery strategies. & Reported parameters are not certified for Manzil's network. & Implement explicit safe-transition and recovery states.\\
Louati et al. \cite{el2018emergency} & Queue-aware multi-agent emergency crossing at signalized intersections. & Its multi-vehicle and multi-agent scope exceeds the core case. & Retain queue awareness and downstream checks.\\
\bottomrule
\end{longtable}
}

\section{Conclusion}

The review supports a conservative technical strategy in which evidence is built in stages. Transparent fixed-time and pressure-inspired controllers must operate correctly before a learned controller is compared with them, and measurement definitions must be fixed before any optimization result is interpreted. Video experiments should assess the complete observation chain rather than quote detector accuracy as a traffic-management result. Simulation experiments should use matched seeds, identical safety constraints, and held-out scenarios so that controller differences are not confused with different demand. Emergency evaluation must also separate the effect of route selection from the effect of signal priority, because combining both changes in one comparison would prevent a clear explanation of any improvement. The cited literature justifies the selected building blocks but does not establish their performance for Manzil's network, cameras, or computing resources. Chapter \ref{ch:srs} therefore converts these findings into traceable requirements, while Chapter \ref{ch:method} defines the measurements and controlled comparisons needed to answer the research questions without claiming field deployment or certified safety evidence.
